\documentclass[10pt]{extarticle}
\usepackage[a4paper,margin=1.5cm]{geometry}
\usepackage[T1]{fontenc}
\usepackage{amsmath,amssymb}
\usepackage{graphicx}
\usepackage[export]{adjustbox}
\usepackage{multicol}
\usepackage{enumitem}
\usepackage{booktabs,array,tabularx}
\setlength{\columnsep}{0.7cm}
\setlength{\parindent}{0pt}
\setlength{\parskip}{3pt}
\pagestyle{empty}
\makeatletter
\newcommand{\qfigdef}[4]{\@namedef{qf@file@#1}{#2}\@namedef{qf@w@#1}{#3}\@namedef{qf@wide@#1}{#4}}
\newcommand{\qf@check}[1]{\@ifundefined{qf@file@#1}{\PackageError{score2}{Unknown figure id #1}{}}{}}
\newcommand{\qf@img}[4][]{\includegraphics[width=#2\dimexpr\@nameuse{qf@w@#4}pt\relax,max width=#3,max height=0.45\textheight,#1]{\@nameuse{qf@file@#4}}}
\newif\ifqf@full
\newif\ifqf@boxed
\newsavebox{\qf@box}
\newcommand{\qf@open}{\global\qf@fulltrue\par\global\qf@boxedtrue
  \setbox\qf@box\vbox\bgroup\hsize\textwidth\linewidth\textwidth}
\newcommand{\qf@emit}[1]{\par\egroup\global\qf@boxedfalse
  \ifdim\dimexpr\ht\qf@box+\dp\qf@box\relax<#1\textheight\noindent\usebox{\qf@box}\par\else\unvbox\qf@box\fi}
\newcommand{\qfullsplit}{\ifqf@boxed\qf@emit{0.85}\fi}
\newcommand{\qf@close}{\ifqf@boxed\qf@emit{0.45}\fi\global\qf@fullfalse}
\newcommand{\qfullbegin}{\end{multicols}\qf@open}
\newcommand{\qfullend}{\qf@close\begin{multicols}{2}}
\newcommand{\qfullstart}{\qf@open}
\newcommand{\qfulldone}{\qf@close}
\newcommand{\qfig}[2][1]{\qf@check{#2}%
  \ifqf@full\qf@blk{#1}{#2}\else
  \ifnum\@nameuse{qf@wide@#2}=1 \expandafter\@firstoftwo\else\expandafter\@secondoftwo\fi
  {\end{multicols}\begin{center}\qf@img{#1}{\textwidth}{#2}\end{center}\begin{multicols}{2}}%
  {\qf@blk{#1}{#2}}\fi}
\newcommand{\qf@blk}[2]{\par\smallskip{\centering\qf@img{#1}{\linewidth}{#2}\par}\smallskip}
\newcommand{\qfigi}[2][1]{\qf@check{#2}\qf@img[valign=c]{#1}{\linewidth}{#2}}
\makeatother
\newcommand{\papersec}[1]{\par\bigskip{\centering\Large\bfseries PAPER-#1\par}\medskip}
\newcommand{\partsec}[2]{\par\medskip\noindent{\large\bfseries PART-#1: #2}\par\smallskip}
\newcommand{\sect}[1]{\par\smallskip\noindent\textbf{#1}\par}
\newcommand{\qn}[1]{\par\smallskip\noindent\textbf{#1.}~}
\newcommand{\paraq}[2]{\par\smallskip\noindent\textbf{Paragraph for Questions #1 and #2}\par}
\newenvironment{qoptions}{\begin{itemize}[leftmargin=1.8em,itemsep=1pt,topsep=2pt,label={}]}{\end{itemize}}
\newcommand{\opt}[2]{\item[(#1)] #2}
% figure registry (generated by score2_tools.py finalize)
\qfigdef{p1_phy_s1_q04_1}{Score-II - Test 22 (14 Mar 2024)-img/p1_phy_s1_q04_1.png}{73.2}{0}
\qfigdef{p1_phy_s1_q07_1}{Score-II - Test 22 (14 Mar 2024)-img/p1_phy_s1_q07_1.png}{84.2}{0}
\qfigdef{p1_phy_s1_q12_1}{Score-II - Test 22 (14 Mar 2024)-img/p1_phy_s1_q12_1.png}{126.7}{0}
\qfigdef{p1_phy_s1_q15_1}{Score-II - Test 22 (14 Mar 2024)-img/p1_phy_s1_q15_1.png}{76.0}{0}
\qfigdef{p1_phy_s1_q16_1}{Score-II - Test 22 (14 Mar 2024)-img/p1_phy_s1_q16_1.png}{63.0}{0}
\qfigdef{p1_phy_s2_q04_1}{Score-II - Test 22 (14 Mar 2024)-img/p1_phy_s2_q04_1.png}{65.7}{0}
\qfigdef{p1_phy_s2_q06_1}{Score-II - Test 22 (14 Mar 2024)-img/p1_phy_s2_q06_1.png}{80.8}{0}
\qfigdef{p1_phy_s2_q10_1}{Score-II - Test 22 (14 Mar 2024)-img/p1_phy_s2_q10_1.png}{50.0}{0}
\qfigdef{p1_chem_s1_q14_2}{Score-II - Test 22 (14 Mar 2024)-img/p1_chem_s1_q14_2.png}{49.9}{0}
\qfigdef{p1_chem_s1_q14_1}{Score-II - Test 22 (14 Mar 2024)-img/p1_chem_s1_q14_1.png}{36.8}{0}
\qfigdef{p1_chem_s1_q14_3}{Score-II - Test 22 (14 Mar 2024)-img/p1_chem_s1_q14_3.png}{49.7}{0}
\qfigdef{p1_chem_s1_q14_4}{Score-II - Test 22 (14 Mar 2024)-img/p1_chem_s1_q14_4.png}{44.5}{0}
\qfigdef{p1_chem_s1_q16_1}{Score-II - Test 22 (14 Mar 2024)-img/p1_chem_s1_q16_1.png}{65.3}{0}
\qfigdef{p1_chem_s1_q16_B}{Score-II - Test 22 (14 Mar 2024)-img/p1_chem_s1_q16_B.png}{75.7}{0}
\qfigdef{p1_chem_s1_q17_1}{Score-II - Test 22 (14 Mar 2024)-img/p1_chem_s1_q17_1.png}{156.8}{0}
\qfigdef{p1_chem_s1_q18_1}{Score-II - Test 22 (14 Mar 2024)-img/p1_chem_s1_q18_1.png}{70.8}{0}
\qfigdef{p1_chem_s1_q19_1}{Score-II - Test 22 (14 Mar 2024)-img/p1_chem_s1_q19_1.png}{117.6}{0}
\qfigdef{p1_chem_s1_q19_2}{Score-II - Test 22 (14 Mar 2024)-img/p1_chem_s1_q19_2.png}{35.7}{0}
\qfigdef{p1_chem_s1_q19_3}{Score-II - Test 22 (14 Mar 2024)-img/p1_chem_s1_q19_3.png}{48.5}{0}
\qfigdef{p1_chem_s1_q19_4}{Score-II - Test 22 (14 Mar 2024)-img/p1_chem_s1_q19_4.png}{35.7}{0}
\qfigdef{p1_chem_s2_q07_1}{Score-II - Test 22 (14 Mar 2024)-img/p1_chem_s2_q07_1.png}{107.2}{0}
\qfigdef{p1_chem_s2_q08_1}{Score-II - Test 22 (14 Mar 2024)-img/p1_chem_s2_q08_1.png}{124.3}{0}
\qfigdef{p1_chem_s2_q09_1}{Score-II - Test 22 (14 Mar 2024)-img/p1_chem_s2_q09_1.png}{68.8}{0}
\qfigdef{p1_chem_s2_q10_1}{Score-II - Test 22 (14 Mar 2024)-img/p1_chem_s2_q10_1.png}{28.3}{0}
\begin{document}
\begin{center}{\Large\bfseries JEE (Main+Advanced) Score-II --- Test 22 (14 Mar 2024)}\end{center}
\begin{multicols}{2}
\partsec{1}{PHYSICS}

\sect{Section-I: Single Correct}

\qn{1} The temperature of 5 mole gas changed from $100^\circ C$ to $120^\circ C$. Change in internal energy is 80 J. The total heat capacity of the gas at constant volume will be :-
\begin{qoptions}
\opt{A}{8 $JK^{-1}$}
\opt{B}{0.8 $JK^{-1}$}
\opt{C}{4 $JK^{-1}$}
\opt{D}{0.4 $JK^{-1}$}
\end{qoptions}

\qn{2} 540 g of ice at $0^\circ C$ is mixed with 540 g water at $80^\circ C$. The final temperature of the mixture is :-
\begin{qoptions}
\opt{A}{$0^\circ C$}
\opt{B}{$40^\circ C$}
\opt{C}{$80^\circ C$}
\opt{D}{less than $0^\circ C$}
\end{qoptions}

\qn{3} Two identical thin rings, each of radius R meter are coaxially placed at distance R meter apart. If $Q_1$ and $Q_2$ coulomb are respectively the charges uniformly spread on the two rings, the work done in moving a charge q from the centre of one ring to that of the other is :-
\begin{qoptions}
\opt{A}{Zero}
\opt{B}{$q(Q_2-Q_1)(\sqrt2-1)/(\sqrt2\cdot4\pi\varepsilon_0R)$}
\opt{C}{$q\sqrt2(Q_1+Q_2)/4\pi\varepsilon_0R$}
\opt{D}{$q(Q_1-Q_2)(\sqrt2+1)/(\sqrt2\cdot4\pi\varepsilon_0R)$}
\end{qoptions}

\qn{4} In this figure the resistance of the coil of galvanometer G is $2\Omega$. The emf of the cell is 4V. The ratio of potential difference across $C_1$ and $C_2$ is:
\qfig{p1_phy_s1_q04_1}
\begin{qoptions}
\opt{A}{1}
\opt{B}{$\dfrac45$}
\opt{C}{$\dfrac34$}
\opt{D}{$\dfrac54$}
\end{qoptions}

\qn{5} A copper wire of length `$\ell$' and radius `r' is nickel plated till its final radius is 2r. If the resistivity of the copper and nickle are $\rho_c$ and $\rho_n$, then find the equivalent resistance of wire between its ends ?
\begin{qoptions}
\opt{A}{$\dfrac{\rho_c\rho_n\ell}{\pi r^2(\rho_n+3\rho_c)}$}
\opt{B}{$\dfrac{\rho_c\rho_n\ell}{\pi r^2(\rho_n+4\rho_c)}$}
\opt{C}{$\dfrac{\rho_c\rho_n\ell}{4\pi r^2(\rho_n+3\rho_c)}$}
\opt{D}{None of these}
\end{qoptions}

\qn{6} A solenoid of 10 henry inductance and 2 ohm resistance, is connected to a 10 volt battery. In how much time the magnetic energy will increases to one forth of the maximum value :--
\begin{qoptions}
\opt{A}{3.5 s}
\opt{B}{2.5 s}
\opt{C}{5.5 s}
\opt{D}{7.5 s}
\end{qoptions}

\qn{7} A wire ABCDEF with each side of length `$\ell$' bent as shown in figure and carrying a current I. If it is placed in a uniform magnetic field B which is parallel to +y direction. Magnetic force experienced by the wire is:-
\qfig{p1_phy_s1_q07_1}
\begin{qoptions}
\opt{A}{$BI\ell$, along $+z$ direction}
\opt{B}{$BI\ell$, along $-z$ direction}
\opt{C}{$2BI\ell$, along $+z$ direction}
\opt{D}{$2BI\ell$, along $-z$ direction}
\end{qoptions}

\qn{8} An L--C--R circuit has L = 10 mH, R = 3 ohm and C = $1\mu F$ connected in series to a source of $V=15\cos\omega t$ volt. The current-amplitude and the average power dissipated per cycle at a frequency 10\% lower than the resonant frequency will respectively be :-
\begin{qoptions}
\opt{A}{0.704A, 0.744W}
\opt{B}{0.704A, 0.704W}
\opt{C}{7.04A, 7.44W}
\opt{D}{70.4A, 74.4W}
\end{qoptions}

\qn{9} A radio wave of frequency 90 MHz (FM) enters a ferrite rod. In ferrite, $\varepsilon_r=10^3$ and $\mu_r=10$, then the velocity and wavelength in ferrite are :-
\begin{qoptions}
\opt{A}{$3\times10^6\ ms^{-1}$, $3.33\times10^{-2}$ m}
\opt{B}{$3\times10^6\ ms^{-1}$, $3.33\times10^{-4}$ m}
\opt{C}{$3\times10^6\ ms^{-1}$, $3.33\times10^{-3}$ m}
\opt{D}{None of the above}
\end{qoptions}

\qn{10} When photons of energy 4.25 eV strike the surface of a metal A, the ejected photoelectrons have maximum kinetic energy $T_A$ eV and de-Broglie wavelength $\lambda_A$. The maximum kinetic energy of photoelectrons liberated from another metal B by photons of energy 4.70 eV is $T_B=(T_A-1.50)$ eV and de-Broglie wavelength is $\lambda_B$. If the de-Broglie wavelengths of these photoelectrons are related as $\lambda_B=2\lambda_A$, then :-\\
(a) The work function of A is 2.25 eV\\
(b) The work function of B is 4.20 eV\\
(c) $T_A=2.00$ eV\\
(d) $T_B=2.75$ eV
\begin{qoptions}
\opt{A}{a, c}
\opt{B}{b, d}
\opt{C}{a, b, c}
\opt{D}{b, c, d}
\end{qoptions}

\qn{11} With what potential an electron should be accelerated so that its de Broglie wave-length becomes equal to the wave length of second line of Balmer series for $He^+$ ion ?
\begin{qoptions}
\opt{A}{$\dfrac{9R^2h^2}{32me}$}
\opt{B}{$\dfrac{162R^2h^2}{25me}$}
\opt{C}{$\dfrac{81R^2h^2}{32me}$}
\opt{D}{$\dfrac{32R^2h^2}{81me}$}
\end{qoptions}

\qn{12} Output for the following Boolean circuit is :
\qfig{p1_phy_s1_q12_1}
\begin{qoptions}
\opt{A}{$\overline{AB}$}
\opt{B}{$\overline{AB}+\overline{AB}$}
\opt{C}{$\overline{A\overline{B}}+\overline{\overline{A}B}$}
\opt{D}{$A+B+\overline{AB}$}
\end{qoptions}

\qn{13} Two beams of light having intensities I and 4I interfere to produce a fringe pattern on a screen. The phase difference between the beam is $\dfrac\pi2$ at point A and $2\pi$ at point B. Then find out the difference between the resultant intensities at A and B.
\begin{qoptions}
\opt{A}{2I}
\opt{B}{5I}
\opt{C}{I}
\opt{D}{4I}
\end{qoptions}

\qn{14} If the length of rod A is $3.25\pm0.01$ cm and that of B is $4.19\pm0.01$ cm, then the rod B is longer than rod A by :-
\begin{qoptions}
\opt{A}{$0.94\pm0.00$ cm}
\opt{B}{$0.94\pm0.01$ cm}
\opt{C}{$0.94\pm0.02$ cm}
\opt{D}{$0.94\pm0.005$ cm}
\end{qoptions}

\qn{15} Figure shows two pulley arrangements for lifting a mass m. In (a) the mass is lifted by attaching a mass 2 m while in (b) the mass is lifted by pulling the other end with a downward force F = 2 mg, If $f_a$ \& $f_b$ are the accelerations of the two masses then (Assume string is massless and pulley is ideal)
\qfig{p1_phy_s1_q15_1}
\begin{qoptions}
\opt{A}{$f_a=f_b$}
\opt{B}{$f_a=f_b/2$}
\opt{C}{$f_a=f_b/3$}
\opt{D}{$f_a=2f_b$}
\end{qoptions}

\qn{16} A block of mass M is sliding down the plane. Coefficient of static friction is $\mu_s$ and kinetic friction is $\mu_k$. Then friction force acting on the block is :-
\qfig{p1_phy_s1_q16_1}
\begin{qoptions}
\opt{A}{$\mu_sMg\cos\theta$}
\opt{B}{$(F+Mg)\sin\theta$}
\opt{C}{$\mu_k(F+Mg)\cos\theta$}
\opt{D}{$(Mg+F)\tan\theta$}
\end{qoptions}

\qn{17} A particle travels such that $\dfrac{dV}{dt}=-a\sqrt v$ where a is constant. Its initial velocity is $v_0$. The time at which it stops is :-
\begin{qoptions}
\opt{A}{$\dfrac{2\sqrt{v_0}}{a}$}
\opt{B}{$\dfrac{\sqrt{v_0}}{a}$}
\opt{C}{$2a\sqrt{v_0}$}
\opt{D}{$a\sqrt{v_0}$}
\end{qoptions}

\qn{18} The centre of mass of a non uniform rod of length L whose mass per unit length $\lambda=\left(\dfrac{Kx^2}{L}\right)$ where K is a constant and x is the distance from one end is :-
\begin{qoptions}
\opt{A}{$\dfrac{3L}{4}$}
\opt{B}{$\dfrac{L}{8}$}
\opt{C}{$\dfrac{L}{K}$}
\opt{D}{$\dfrac{3L}{K}$}
\end{qoptions}

\qn{19} A circular disc of radius R and thickness $\dfrac R6$ has moment of inertia I about an axis passing through its centre and perpendicular to its plane. It is melted and recasted into a solid sphere. The moment of inertia of the sphere about its diameter as axis of rotation is:-
\begin{qoptions}
\opt{A}{I}
\opt{B}{$\dfrac{2I}{8}$}
\opt{C}{$\dfrac{I}{5}$}
\opt{D}{$\dfrac{I}{10}$}
\end{qoptions}

\qn{20} A satellite is orbiting the earth in a circular orbit of radius r. Its :-\\
(a) Kinetic energy varies as $\dfrac1r$\\
(b) Angular momentum varies as $\dfrac{1}{\sqrt r}$\\
(c) Linear momentum varies as $\dfrac{1}{\sqrt r}$\\
(d) Frequency of revolution varies as $\dfrac{1}{r^{3/2}}$
\begin{qoptions}
\opt{A}{Only d}
\opt{B}{a, c, d}
\opt{C}{Only a, b}
\opt{D}{Only b}
\end{qoptions}

\sect{Section-II: Numerical Value}

\qn{1} If the volume of a block of metal changes by 0.12\% when it is heated through $20^\circ C$, the coefficient of linear expansion (in $^\circ C^{-1}$) of the metal is $x\times10^{-5}$, find the value of x ?

\qn{2} The volume of a metal sphere increases by 0.24\% when its temperature is raised by $40^\circ C$. The coefficient of linear expansion of the metal is $x\times10^{-5}/^\circ C$, find the value of x ?

\qn{3} The electric field in a region of space is given by $E=(5\hat i+2\hat j)$ N/C. The electric flux through an area of $2m^2$ lying in the YZ plane, in S.I. units is :

\qn{4} The equivalent capacitance between terminals `A' and `B' is $\dfrac{\alpha}{\beta}\dfrac{\varepsilon_0A}{d}$, $\alpha$ \& $\beta$ being least integer. Determine the value of $\alpha+\beta$. (letters have their usual meaning)
\qfig{p1_phy_s2_q04_1}

\qn{5} Mass defect in a $_{11}Na^{23}$ atom is 0.21 amu. Then the binding energy per nucleon $8.505\times10^x$ eV than `x' is :

\qn{6} A zener diode, having breakdown voltage equal to 15 V, is used in a voltage regulator circuit shown in figure. The current through the zener diode is (in mA) :
\qfig{p1_phy_s2_q06_1}

\qn{7} Light travel in two medium $M_1$ \& $M_2$ with speed $1.5\times10^8$ m/s and $2\times10^8$ m/s respectively. The critical angle between them is $\tan^{-1}\left(\dfrac{x}{\sqrt7}\right)$. Find x.

\qn{8} A tuning fork gives 4 beats with 50cm length of a sonometer wire. If the length of the wire is shortened by 1cm, the number of beats is still the same. The frequency of the fork (in Hz) is :-

\qn{9} Two simple pendulums of length 5 m and 20 m respectively are given small linear displacement in one direction at the same time. They will again be in phase when the pendulum of shorter length has completed......... oscillations.

\qn{10} A U-tube in which the cross-sectional area of the limb on the left is one quarter, the limb on the right contains mercury (density 13.6 $g/cm^3$). The level of mercury in the narrow limb is at a distance of 36 cm from the upper end of the tube. What will be the rise in the level of mercury in the right limb if the left limb is filled to the top with water (in cm)
\qfig{p1_phy_s2_q10_1}

\partsec{2}{CHEMISTRY}

\sect{Section-I: Single Correct}

\qn{1} At constant temperature, the equilibrium constant ($K_P$) for the decomposition reaction $N_2O_4\rightleftharpoons2NO_2$ is expressed by $K_P=4x^2P/(1-x^2)$ where P is pressure, x is extent of decomposition. Which of the following statement is true ?
\begin{qoptions}
\opt{A}{$K_P$ increases with increase of P}
\opt{B}{$K_P$ increases with increase of x}
\opt{C}{$K_P$ increases with decrease of x}
\opt{D}{$K_P$ remains constant with change in P or x}
\end{qoptions}

\qn{2} Which of the following is true statement ?
\begin{qoptions}
\opt{A}{Rate of reaction does not changes on changing stoichiometry of reaction}
\opt{B}{Rate can be negative}
\opt{C}{If rate of appearance of $NH_3$ is 1.7 kg/hr in the reaction $N_2+H_2\longrightarrow NH_3$ then rate of disappearance of $H_2$ will be 42.5 gm/min}
\opt{D}{Rate of disappearance of reactants or rate of appearance of products does not changes on changing stoichiometry of the reaction}
\end{qoptions}

\qn{3} In 500 ml of 0.2M-NaCl solution, 1500 ml of 0.4M $MgCl_2$ solution is mixed. The only incorrect final concentration of ion is (assume volume to be additive) :
\begin{qoptions}
\opt{A}{$[Na^+]=0.2$ M}
\opt{B}{$[Mg^{2+}]=0.3$ M}
\opt{C}{$[Cl^-]=0.65$ M}
\opt{D}{$[Mg^{2+}]=7.2$ g/L}
\end{qoptions}

\qn{4} The solubility of $PbF_2$ (formula weight = 245) is 0.46g/L. What is the solubility product :--
\begin{qoptions}
\opt{A}{$1.1\times10^{-10}$}
\opt{B}{$2.6\times10^{-8}$}
\opt{C}{$1.1\times10^{-7}$}
\opt{D}{$6.8\times10^{-9}$}
\end{qoptions}

\qn{5} The vapour pressure of $CS_2$ at $50^\circ C$ is 854 torr and a solution of 2.0 gm sulphur in 100 gm of $CS_2$ has vapour pressure 848.9 torr. If the formula of sulphur molecule is $S_n$, then calculate the value of "n" (approx) (At. mass of S = 32) :-\\
\relax[Assume very dilute solution]
\begin{qoptions}
\opt{A}{2}
\opt{B}{4}
\opt{C}{8}
\opt{D}{6}
\end{qoptions}

\qn{6} 18 gm glucose is completely combusted in bomb calorimeter, of heat capacity 1400 kJ/K, temperature changes from $27^\circ C$ to $27.2\ ^\circ C$ magnitude of standard enthalpy of combustion of glucose in kJ/mol. [R = 8.314 J/mol-K]
\begin{qoptions}
\opt{A}{1200}
\opt{B}{2000}
\opt{C}{2800}
\opt{D}{1400}
\end{qoptions}

\qn{7} To which of the following atom, the attachment of electron is most difficult ?
\begin{qoptions}
\opt{A}{Radon}
\opt{B}{Nitrogen}
\opt{C}{Oxygen}
\opt{D}{Radium}
\end{qoptions}

\qn{8} Which of the following statement is not correct ?
\begin{qoptions}
\opt{A}{$Lu^{+3}$ has the strongest tendency toward complex formation among trivalent lanthanoid ions}
\opt{B}{Ce has maximum composition in misch metal}
\opt{C}{f-block elements can have electrons from $f^0$ to $f^{14}$}
\opt{D}{Nd, Np and Nb all are f-block elements}
\end{qoptions}

\qn{9} Why is HCl not used to make the medium acidic in oxidation reactions of $KMnO_4$ in acidic medium?
\begin{qoptions}
\opt{A}{Both HCl and $KMnO_4$ act as oxidising agents.}
\opt{B}{$KMnO_4$ oxidises HCl into $Cl_2$ which is also an oxidising agent.}
\opt{C}{$KMnO_4$ is a weaker oxidising agent than HCl.}
\opt{D}{$KMnO_4$ acts as a reducing agent in the presence of HCl.}
\end{qoptions}

\qn{10} The correct order of the O--O bond length in $O_2$, $H_2O_2$ and $O_3$ is :-
\begin{qoptions}
\opt{A}{$O_2>O_3>H_2O_2$}
\opt{B}{$O_3>H_2O_2>O_2$}
\opt{C}{$H_2O_2>O_3>O_2$}
\opt{D}{$H_2O_2>O_2>O_3$}
\end{qoptions}

\qn{11} Which of the following complex have $t_{2g}^4,e_g^2$ configuration ?
\begin{qoptions}
\opt{A}{$[CoF_6]^{-3}$}
\opt{B}{$[Co(H_2O)_6]^{+3}$}
\opt{C}{$[Co(C_2O_4)_3]^{-3}$}
\opt{D}{$[Co(CN)_6]^{-3}$}
\end{qoptions}

\qn{12} Which of the following complex will show optical activity ?
\begin{qoptions}
\opt{A}{trans-$[Co(NH_3)_4Cl_2]^+$}
\opt{B}{$[Cr(H_2O)_6]^{3+}$}
\opt{C}{Cis-$[Co(NH_3)_2(en)_2]^{3+}$}
\opt{D}{trans-$[Co(NH_3)_2(en)_2]^{3+}$}
\end{qoptions}

\qn{13} The orbital angular momentum of electron in `d' orbital is equal to :-
\begin{qoptions}
\opt{A}{$\sqrt2\hbar$}
\opt{B}{$2\sqrt3\hbar$}
\opt{C}{$0\hbar$}
\opt{D}{$\sqrt6\hbar$}
\end{qoptions}

\qn{14} In which molecule, cleavage by $HIO_4$ is not observed ?
\begin{qoptions}
\opt{A}{\qfigi{p1_chem_s1_q14_2}}
\opt{B}{\qfigi{p1_chem_s1_q14_1}}
\opt{C}{\qfigi{p1_chem_s1_q14_3}}
\opt{D}{\qfigi{p1_chem_s1_q14_4}}
\end{qoptions}

\qn{15} Match the following reactants in Column I with the corresponding reaction products in Column II :-
\begin{center}
\begin{tabular}{|c|p{3.3cm}|c|p{2.1cm}|}
\hline
\multicolumn{2}{|c|}{Column-I} & \multicolumn{2}{c|}{Column-II} \\
\hline
(a) & Benzene $+Cl_2\xrightarrow{AlCl_3}$ & (1) & Benzoic acid \\
\hline
(b) & Benzene $+CH_3Cl\xrightarrow{AlCl_3}$ & (2) & Methyl Phenyl ketone \\
\hline
(c) & Benzene $+CH_3COCl$\newline\resizebox{1.6cm}{!}{$\xrightarrow{AlCl_3}$} & (3) & Toluene \\
\hline
(d) & Toluene\newline\resizebox{1.6cm}{!}{$\xrightarrow{KMnO_4/NaOH}$} & (4) & Chloro\-benzene \\
\hline
 & & (5) & Benzene hexachloride \\
\hline
\end{tabular}
\end{center}
Codes\\
\begin{tabular}{ccccc}
 & a & b & c & d\\
\end{tabular}
\begin{qoptions}
\opt{A}{4\quad3\quad2\quad1}
\opt{B}{1\quad2\quad3\quad4}
\opt{C}{4\quad3\quad1\quad5}
\opt{D}{1\quad5\quad2\quad3}
\end{qoptions}

\qn{16} Which of the following is not correctly matched :-
\begin{qoptions}
\opt{A}{\qfigi{p1_chem_s1_q16_1} Reimer-Tiemann formylation}
\opt{B}{\qfigi{p1_chem_s1_q16_B} Gatterman Koch reaction}
\opt{C}{$R-N\equiv C\xrightarrow{CHCl_3+KOH}$ Isocyanide test}
\opt{D}{$R-N\equiv C\xrightarrow[\text{(ii)}H_3O^+]{\text{(i)}SnCl_2+HCl}$ Mendius Reaction}
\end{qoptions}

\qn{17} \qfig{p1_chem_s1_q17_1}
\begin{qoptions}
\opt{A}{P$\to$1;Q$\to$2;R$\to$3;S$\to$4}
\opt{B}{P$\to$4;Q$\to$3;R$\to$2;S$\to$1}
\opt{C}{P$\to$2;Q$\to$4;R$\to$3;S$\to$1}
\opt{D}{P$\to$2;Q$\to$3;R$\to$4;S$\to$1}
\end{qoptions}

\qn{18} The given osazone can be obtained by :-
\qfig{p1_chem_s1_q18_1}
\begin{qoptions}
\opt{A}{D--glucose}
\opt{B}{D--mannose}
\opt{C}{D--Idose}
\opt{D}{Both (A) \& (B)}
\end{qoptions}

\qn{19} \qfig{p1_chem_s1_q19_1}
Compound (B) is :-
\begin{qoptions}
\opt{A}{\qfigi{p1_chem_s1_q19_2}}
\opt{B}{\qfigi{p1_chem_s1_q19_3}}
\opt{C}{\qfigi{p1_chem_s1_q19_4}}
\opt{D}{(A) and (C) both}
\end{qoptions}

\qn{20} $CH_3CH=CHCHO\xrightarrow[\text{Aldol condensation}]{OH^-}\xrightarrow{\Delta}$;\\
Product A is :-
\begin{qoptions}
\opt{A}{$CH_3(CH=CH)_3CHO$}
\opt{B}{$CH_3CH_2CH_2(CH=CH)_2CHO$}
\opt{C}{$CH_3(CH_2CH_2)_3CH=CH-CHO$}
\opt{D}{None is correct}
\end{qoptions}

\sect{Section-II: Numerical Value}

\qn{1} Calculate value of $\lambda_m^0\ (\Omega^{-1}cm^2mole^{-1})$ for $CaCl_2$ from following data :
\begin{center}
\begin{tabular}{|c|c|c|}
\hline
Conc. (mole/litre) & 0.25 & 1 \\
\hline
$\lambda_m\ (\Omega^{-1}cm^2mole^{-1})$ & 300 & 280 \\
\hline
\end{tabular}
\end{center}

\qn{2} 3 moles of $FeSO_4$ are oxidised by "a" mole of $K_2Cr_2O_7$ in acidic medium whereas 3 moles of $FeC_2O_4$ are oxidised by "b" mole of $K_2Cr_2O_7$ in acidic medium, the ratio of "b" and "a" is :-

\qn{3} A quantity of aluminium has a mass of 54 gm. What is the mass of same number of magnesium atoms? (Given : At. wt. of Mg = 24)

\qn{4} For which orbit number of $Be^{3+}$, value of radius is equal to 0.529 \AA\ ?

\qn{5} Count basic hydroxides NaOH, ClOH, KOH, $Ca(OH)_2$, $Zn(OH)_2$, $(OH)_2SO_2$, $HO(NO_2)$, $Al(OH)_3$

\qn{6} Find the total number having both $p_\pi-p_\pi$ \& $p_\pi-d_\pi$ type of bond.\\
$SO_3$, $SO_2$, $SO_4^{2-}$, $SO_3^{2-}$, $XeO_3$, $XeO_4$, $POCl_3$, $SNF_3$

\qn{7} In the following reaction
\qfig{p1_chem_s2_q07_1}
The number of $\alpha$-hydrogen present in B compound.

\qn{8} How many of the following compounds is are aromatic ?
\qfig{p1_chem_s2_q08_1}

\qn{9} How many moles of RMgX consumed in the given compound ?
\qfig{p1_chem_s2_q09_1}

\qn{10} \qfig{p1_chem_s2_q10_1}
Total number of geometrical isomers of above compound are

\partsec{3}{MATHEMATICS}

\sect{Section-I: Single Correct}

\qn{1} If $a+b+c=3$ and $a>0$, $b>0$, $c>0$ then the greatest value of $a^2b^2c^2$ is :-
\begin{qoptions}
\opt{A}{$\dfrac{3^{10}2^4}{7^7}$}
\opt{B}{1}
\opt{C}{$\dfrac{3^82^4}{7^7}$}
\opt{D}{None of these}
\end{qoptions}

\qn{2} If $\alpha,\beta$ are roots of eq. $375x^2-25x-2=0$ \& $S_n=\alpha^n+\beta^n$ then the value of $\displaystyle\lim_{n\to\infty}\sum_{r=1}^nS_r=$
\begin{qoptions}
\opt{A}{$\dfrac34$}
\opt{B}{$\dfrac{1}{12}$}
\opt{C}{0}
\opt{D}{15}
\end{qoptions}

\qn{3} If A is matrix of order $(3\times3)$ and $|A|=2$ then value of $|adj\ adj\ adj\ adj\ adj\ (2A)|$ is equal to : --
\begin{qoptions}
\opt{A}{$2^9$}
\opt{B}{$2^{128}$}
\opt{C}{$2^{32}$}
\opt{D}{None}
\end{qoptions}

\qn{4} Let 13 seats in a straight line are to be occupied by 4 persons. The number of ways in which there should be exactly four empty seats and no other person between two particular persons is-
\begin{qoptions}
\opt{A}{$15(^6C_3)$}
\opt{B}{$30(^6C_3)$}
\opt{C}{$16(^7C_2)$}
\opt{D}{$32(^7C_2)$}
\end{qoptions}

\qn{5} In a pack of 52 cards one by one cards are randomly drawn untill all red court cards are not drawn, then find the probability that this process is completed in 17th draw
\begin{qoptions}
\opt{A}{$\dfrac{^{12}C_6\cdot^{40}C_{11}}{^{52}C_{16}\cdot^{36}C_{17}}$}
\opt{B}{$\dfrac{1}{46}\dfrac{^6C_5\cdot^{46}C_{11}}{^{52}C_{16}}$}
\opt{C}{$\dfrac{^{46}C_{11}}{^{52}C_{16}\cdot6}$}
\opt{D}{None}
\end{qoptions}

\qn{6} The value of $i^{1+3+5+\ldots+(2n+1)}$ is :-
\begin{qoptions}
\opt{A}{i if n is even, $-i$ if n is odd}
\opt{B}{1 if n is even, $-1$ if n is odd}
\opt{C}{1 if n is odd, $-1$ if n is even}
\opt{D}{i if n is even, 1 if n is odd}
\end{qoptions}

\qn{7} If the maximum value of $y=\dfrac{7+6\tan x-\tan^2x}{(1+\tan^2x)}$ is $\lambda$ then the value of $\log_{\sqrt2}(\lambda)$ is :
\begin{qoptions}
\opt{A}{0}
\opt{B}{6}
\opt{C}{8}
\opt{D}{1}
\end{qoptions}

\qn{8} S.D. of $x_1,x_2,\ldots,x_n$ is 2, if $\displaystyle\sum_{i=1}^nx_i=20$ and $\displaystyle\sum_{i=1}^nx_i^2=100$. Then value of n is
\begin{qoptions}
\opt{A}{5 or 15}
\opt{B}{10 or 20}
\opt{C}{5 or 10}
\opt{D}{5 or 20}
\end{qoptions}

\qn{9} Let R be the relation on the set R of all real numbers defined by a R b iff $|a-b|\le1$. Then R is :-
\begin{qoptions}
\opt{A}{Reflexive and symmetric}
\opt{B}{Symmetric only}
\opt{C}{Transitive only}
\opt{D}{Anti symmetric only}
\end{qoptions}

\qn{10} Let $f(x)=2\cos^{-1}x+4\cot^{-1}x-3x^2-2x+10$, $x\in[-1,1]$. If $[a,b]$ is the range of the function then $4a-b$ is equal to:
\begin{qoptions}
\opt{A}{11}
\opt{B}{$11-\pi$}
\opt{C}{$11+\pi$}
\opt{D}{$15-\pi$}
\end{qoptions}

\qn{11} $\sin^{-1}\left(\sin\dfrac{2\pi}{3}\right)+\cos^{-1}\left(\cos\dfrac{7\pi}{6}\right)+\tan^{-1}\left(\tan\dfrac{3\pi}{4}\right)$ is equal to :
\begin{qoptions}
\opt{A}{$\dfrac{11\pi}{12}$}
\opt{B}{$\dfrac{17\pi}{12}$}
\opt{C}{$\dfrac{31\pi}{12}$}
\opt{D}{$-\dfrac{3\pi}{4}$}
\end{qoptions}

\qn{12} Let $f(x)=\begin{cases}x^3-x^2+10x-7, & x\le1\\-2x+\log_2(b^2-4), & x>1\end{cases}$\\
Then the set of all values of b, for which f(x) has maximum value at $x=1$, is :
\begin{qoptions}
\opt{A}{$(-6,-2)$}
\opt{B}{$(2,6)$}
\opt{C}{$[-6,-2)\cup(2,6]$}
\opt{D}{$[-\sqrt6,-2)\cup(2,\sqrt6]$}
\end{qoptions}

\qn{13} Let $\lambda^*$ be the largest value of $\lambda$ for which the function $f_\lambda(x)=4\lambda x^3-36\lambda x^2+36x+48$ is increasing for all $x\in R$. Then $f_{\lambda^*}(1)+f_{\lambda^*}(-1)$ is equal to :
\begin{qoptions}
\opt{A}{36}
\opt{B}{48}
\opt{C}{64}
\opt{D}{72}
\end{qoptions}

\qn{14} The area of the region $S=\{(x,y):y^2\le8x,y\ge\sqrt2x,x\ge1\}$ is :
\begin{qoptions}
\opt{A}{$\dfrac{13\sqrt2}{6}$}
\opt{B}{$\dfrac{11\sqrt2}{6}$}
\opt{C}{$\dfrac{5\sqrt2}{6}$}
\opt{D}{$\dfrac{19\sqrt2}{6}$}
\end{qoptions}

\qn{15} Let the solution curve of the differential equation $xdy=\left(\sqrt{x^2+y^2}+y\right)dx$, $x>0$, intersect the line $x=1$ at $y=0$ and the line $x=2$ at $y=\alpha$. Then the value of $\alpha$ is :
\begin{qoptions}
\opt{A}{$\dfrac12$}
\opt{B}{$\dfrac32$}
\opt{C}{$-\dfrac32$}
\opt{D}{$\dfrac52$}
\end{qoptions}

\qn{16} For $I(x)=\displaystyle\int\dfrac{\sec^2x-2022}{\sin^{2022}x}dx$, if $I\left(\dfrac\pi4\right)=2^{1011}$, then
\begin{qoptions}
\opt{A}{$3^{1010}I\left(\dfrac\pi3\right)-I\left(\dfrac\pi6\right)=0$}
\opt{B}{$3^{1010}I\left(\dfrac\pi6\right)-I\left(\dfrac\pi3\right)=0$}
\opt{C}{$3^{1011}I\left(\dfrac\pi3\right)-I\left(\dfrac\pi6\right)=0$}
\opt{D}{$3^{1011}I\left(\dfrac\pi6\right)-I\left(\dfrac\pi3\right)=0$}
\end{qoptions}

\qn{17} Let $\vec a=\hat i+\alpha\hat j+\beta\hat k$, $\alpha,\beta\in R$. Let a vector $\vec b$ be such that the angle between $\vec a$ and $\vec b$ is $\dfrac\pi4$ and $|\vec b|^2=6$, If $\vec a\cdot\vec b=3\sqrt2$, then the value of $(\alpha^2+\beta^2)|\vec a\times\vec b|^2$ is equal to
\begin{qoptions}
\opt{A}{90}
\opt{B}{75}
\opt{C}{95}
\opt{D}{85}
\end{qoptions}

\qn{18} For $0<\theta<\pi/2$, if the eccentricity of the hyperbola $x^2-y^2\text{cosec}^2\theta=5$ is $\sqrt7$ times eccentricity of the ellipse $x^2\text{cosec}^2\theta+y^2=5$, then the value of $\theta$ is :
\begin{qoptions}
\opt{A}{$\dfrac\pi6$}
\opt{B}{$\dfrac{5\pi}{12}$}
\opt{C}{$\dfrac\pi3$}
\opt{D}{$\dfrac\pi4$}
\end{qoptions}

\qn{19} Let $C:x^2+y^2=4$ and $C':x^2+y^2-4\lambda x+9=0$ be two circles. If the set of all values of $\lambda$ so that the circles C and C' intersect at two distinct points, is $R-[a,b]$, then the point $(8a+12,16b-20)$ lies on the curve :
\begin{qoptions}
\opt{A}{$x^2+2y^2-5x+6y=3$}
\opt{B}{$5x^2-y=-11$}
\opt{C}{$x^2-4y^2=7$}
\opt{D}{$6x^2+y^2=42$}
\end{qoptions}

\qn{20} If the shortest distance of the parabola $y^2=4x$ from the centre of the circle $x^2+y^2-4x-16y+64=0$ is d, then $d^2$ is equal to :
\begin{qoptions}
\opt{A}{16}
\opt{B}{24}
\opt{C}{20}
\opt{D}{36}
\end{qoptions}

\sect{Section-II: Numerical Value}

\qn{1} The number of distinct real solution of\\
\resizebox{\linewidth}{!}{$\begin{vmatrix}\sin x&\cos x&\cos x\\\cos x&\sin x&\cos x\\\cos x&\cos x&\sin x\end{vmatrix}=0$}\\
in the interval $-\dfrac\pi4\le x\le\dfrac\pi4$ is-

\qn{2} The number of rational term in expansion of $\left(5^{\frac12}+7^{\frac16}\right)^{65}$ is :-

\qn{3} The value of\\
\resizebox{\linewidth}{!}{$\dfrac{\left(1+\sqrt3\tan1^\circ\right)\left(1+\sqrt3\tan2^\circ\right)(\tan1^\circ+\tan59^\circ)(\tan2^\circ+\tan58^\circ)}{(1+\tan^21^\circ)(1+\tan^22^\circ)}$}\\
is :-

\qn{4} The sum of the maximum and minimum values of the function $f(x)=|5x-7|+[x^2+2x]$ is the interval $\left[\dfrac54,2\right]$, where $[t]$ is the greatest integer $\le t$ is \_\_\_\_\_

\qn{5} Let $f(x)=\left|(x-1)(x^2-2x-3)\right|+x-3$, $x\in R$. If m and M are respectively the number of points of local minimum and local maximum of f in the interval $(0,4)$, then $m+M$ is equal to \_\_\_\_\_.

\qn{6} The value of the integral\\
\resizebox{\linewidth}{!}{$\dfrac{48}{\pi^4}\displaystyle\int_0^\pi\left(\dfrac{3\pi x^2}{2}-x^3\right)\dfrac{\sin x}{1+\cos^2x}dx$}\\
is equal to \_\_\_\_\_\_\_.

\qn{7} If $n(2n+1)\displaystyle\int_0^1(1-x^n)^{2n}dx=1177\displaystyle\int_0^1(1-x^n)^{2n+1}dx$, then $n\in N$ is equal to \_\_\_\_

\qn{8} Let $\displaystyle\text{Max}_{0\le x\le2}\left\{\dfrac{9-x^2}{5-x}\right\}=\alpha$ and $\displaystyle\text{Min}_{0\le x\le2}\left\{\dfrac{9-x^2}{5-x}\right\}=\beta$\\
If $\displaystyle\int_{\beta-\frac83}^{2\alpha-1}\text{Max}\left\{\dfrac{9-x^2}{5-x},x\right\}dx=\alpha_1+\alpha_2\log_e\left(\dfrac{8}{15}\right)$\\
then $\alpha_1+\alpha_2$ is equal to \_\_\_\_\_

\qn{9} A line passes through $A(4,-6,-2)$ and $B(16,-2,4)$. The point $P(a,b,c)$ where a, b, c are non-negative integers, on the line AB lies at a distance of 21 units, from the point A. The distance between the points $P(a,b,c)$ and $Q(4,-12,3)$ is equal to \_\_\_\_.

\qn{10} Let $A(-2,-1)$, $B(1,0)$, $C(\alpha,\beta)$ and $D(\gamma,\delta)$ be the vertices of a parallelogram ABCD. If the point C lies on $2x-y=5$ and the point D lies on $3x-2y=6$, then the value of $|\alpha+\beta+\gamma+\delta|$ is equal to \_\_\_\_\_.

\end{multicols}
\end{document}
